% Imporditud: tools/import_eksperimendid.py
% Allikas: Eesti füüsikaolümpiaadi eksperimendi ülesannete
%   kogu (Rael Kalda, 2023), ülesanne Ü125, lahendus L125.
\setAuthor{Erkki Tempel}
\setRound{lõppvoor}
\setYear{2022}
\setNumber{P E1}
\setDifficulty{3}
\setTopic{geomeetriline optika}
\setVahendid{Nõguslääts, millimeeterpaber (A4), laserpointer}
\setPunktid{12}
\setAllikas{Eksperimendi kogu Ü125, lahendus lk 98}
\setLahendusseis{toores}

\prob{Nõguslääts}
Leidke nõgusläätse fookuskaugus ja kumerusraadius.

\hint

\solu
Asetame läätse millimeeterpaberile. Fookuskauguse mõõtmiseks suuname läätse optilise peateljega paralleelse kiire läätseni ning joonistame millimeeterpaberile murdunud kiire. Kui me suuname läätsele mitu optilise peateljega paralleelset kiirt, siis murdunud kiirte pikendused (joonistame need millimeeterpaberile) kohtuvad fookuses F . Mõõdame millimeeterpaberi abil fookuse F kauguse läätse keskpunktist, mis ongi nõgusläätse fookuskaugus fl.

Kumerusraadiuse leidmiseks talitame samamoodi, kuid märgime millimeeterpaberile läätselt (mis toimub kui nõguspeegel) peegeldunud kiired. Nõguspeegli korral peegli optilise peateljega paralleelsed kiired koonduvad fookuses F . Teades nõguspeegli fookuse F asukohta, saame mõõta fookuskauguse fp. Teades, et kumerusraadius R = 2fp, saame arvutada kumerusraadiuse R.
\probend
